\chapter{Literature Review}
\label{ch:litreview}

\section{Introduction}

This chapter reviews the technology and prior work that shaped the design decisions behind the attendance system. It starts with the fundamentals of RFID, moves through attendance systems that other students and researchers have already built with similar hardware, and then looks specifically at the two design choices that set this project apart from a typical class project: treating the network as unreliable by default, and separating "no Wi-Fi" from "Wi-Fi fine, server down" as two distinct failure states. The chapter closes by identifying the gap in the existing work that this project tries to fill.

\section{Fundamentals of RFID Technology}

Radio Frequency Identification (RFID) is not a new idea. \textcite{landt2005history} traces its roots back to radar research in the 1940s, long before anyone thought to attach a chip to a plastic card. What matters for this project is the more recent, practical form of the technology: a small passive tag containing an antenna and a chip, which draws its power from the electromagnetic field generated by a nearby reader and responds with a unique identifier. \textcite{want2006introduction} describes this exchange clearly — the tag needs no battery of its own, which is exactly why RFID cards can sit in a wallet for years and still work the first time they are tapped. \textcite{finkenzeller2010rfid} goes further into the physics and standards involved, and is a useful reference for anyone who needs to understand why a 13.56 MHz reader such as the MFRC522 used in this project reads a card reliably from a few centimetres away but not from across a room — the read range is a direct consequence of how the tag is powered, not a limitation that can be tuned away in software.

For an attendance system, the property that matters most is that the tag's identifier is fixed and unique. A card cannot be "typed in" wrong, and unlike a signature or a manually entered employee number, it does not depend on the person at the desk transcribing anything correctly. That single property is the foundation the rest of this project's design rests on.

\section{RFID-Based Attendance Systems: Related Work}

RFID attendance systems built around low-cost microcontrollers are now common enough in student and small-scale projects that a reasonably clear pattern has emerged. \textcite{singh2022iot} built an IoT-based digital attendance system using an ESP32 and an RFID reader, logging taps to a live database and reporting that it reduced the manual effort involved in daily attendance capture considerably compared to a paper register. \textcite{yadav2025rfid} took a lighter-weight approach, using an ESP32 and RC522 reader to log attendance directly into a Google Sheet, which keeps the system cheap to deploy but ties every tap to an active internet connection at the moment it happens. \textcite{hutabarat2025implementation} combined the same ESP32-plus-RC522 pairing with a MySQL database and a PHP web interface for a university campus, and specifically credits the system with reducing proxy attendance — a version of the same buddy-punching problem this project is trying to solve for a workplace rather than a classroom. \textcite{turpo2025edurfid} pushed the cost angle furthest, reporting a 94\% cost reduction compared to commercial RFID attendance products by building on a Raspberry Pi and Arduino combination for rural schools in Peru, with a claimed sub-second response time per tap.

Looking across this body of work as a whole, \textcite{ishaq2023iot} carried out a systematic literature review of IoT-based RFID attendance systems and found that most published implementations follow the same broad shape: a microcontroller, an RFID reader, and a live connection to a server or spreadsheet. Very few of the systems reviewed treated connectivity loss as something the design needed to plan for, rather than an edge case to be ignored. That observation lines up with what this project encountered directly during the attachment: a factory entrance is not a lecture hall with dependable Wi-Fi, and a system that only works while the network is up is not good enough for a real workplace.

On the payroll side specifically, \textcite{shukla2019conceptual} proposed a conceptual framework for combining RFID with biometric verification to improve both payroll accuracy and attendance monitoring, arguing that the two technologies address different weaknesses — RFID is fast and cheap, biometrics closes the "borrowed card" gap that RFID alone cannot. That framework is a useful reference point for this project's own limitations, discussed later in Chapter 5, since the system built here uses RFID alone and therefore inherits the same borrowed-card weakness that Shukla and Bhandari set out to solve.

\section{RFID Versus Biometric and Other Identification Methods}

RFID is not the only option for automated attendance capture, and it is worth being honest about where it is weaker than the alternatives. Biometric methods — fingerprint or facial recognition — identify the person directly rather than an object the person is carrying, which closes the door on one employee tapping in for another \parencite{nialabs2025biometric}. \textcite{peniel2025rfid} frames the choice as a trade-off rather than a simple case of one technology being better: RFID is faster per transaction, tolerant of dirty or gloved hands, and considerably cheaper to deploy at scale, while biometric readers are slower, more sensitive to environmental conditions such as moisture or wear, but harder to defeat by simply handing a colleague an object.

For a workplace like Infor Relay Systems (Pvt) Ltd, where dozens of staff may need to clock in within a narrow window at shift change, the speed and low cost of RFID make it the more practical starting point. The trade-off is accepted deliberately here, and flagged clearly in Chapter 5 as an area where a future iteration of the system could add a biometric check for higher-risk cases, rather than treating RFID and biometrics as mutually exclusive choices.

\section{Manual Attendance, Time Theft, and the Case for Automation}

The specific failure mode that motivated this project — a paper register that cannot be trusted — is well documented outside the RFID literature too. \textcite{oloid2025buddy} identifies manual and shared-login attendance systems as the easiest to abuse precisely because they include no mechanism to confirm that the person recording an entry is the person the entry is attributed to. \textcite{qandle2025buddy} describes the resulting practice, buddy punching, as a direct and quantifiable form of time theft: hours paid for that were never actually worked. \textcite{cpcon2025rfid} makes the connection to RFID explicit, arguing that tagging each employee with a unique, hard-to-share credential is one of the more cost-effective ways to close that gap without introducing the friction of a biometric scan at every entrance. This is precisely the problem Infor Relay Systems (Pvt) Ltd faced with its paper register, and it is the reason the objectives in Chapter 1 focus on automation and labour-costing accuracy rather than on any more exotic feature.

\section{Offline-First Design for Embedded and IoT Devices}

The design decision that most separates this project from the related work reviewed above is its treatment of network connectivity. Most published student RFID-attendance projects assume the device is online and write directly to a remote database or spreadsheet on every tap \parencite{yadav2025rfid, hutabarat2025implementation}. That assumption is convenient to build against, but it is fragile: if the Wi-Fi drops for even a few seconds at the exact moment someone taps their card, the tap is lost with no record that it ever happened.

The offline-first pattern used in this project instead treats the network as an optimisation rather than a dependency. \textcite{microsoft2026offline} describes this approach in the context of Azure IoT Edge devices: local storage and a message queue let an edge device keep operating indefinitely while disconnected, syncing everything it collected the moment connectivity returns. \textcite{topic2026offline} summarises the underlying philosophy succinctly — an offline-first system should be fully functional without a network connection and simply \textit{better} once one is available, rather than broken without one. \textcite{metadesk2026offline} makes the same case specifically for IoT hardware, noting that battery- and flash-based local caching is what allows a device deployed in the field to survive router reboots, ISP outages, and intermittent coverage without losing data.

This project applies exactly that pattern at the reader station: every tap is written to the device's local flash storage first, and forwarding it to the server is attempted only after that write succeeds. A queue of unsent taps persists across power cycles and drains automatically once the connection is restored, so a network outage becomes an inconvenience rather than a source of missing attendance records.

\section{Local Service Discovery and Web Protocols}

A reader station that always tries to reach the same fixed IP address breaks the moment a router hands out a different address after a reboot — a problem familiar to anyone who has managed a small office network with mixed hotspot and router use. Multicast DNS, standardised in \textcite{cheshire2013mdns} and its companion service-discovery specification \parencite{cheshire2013dnssd}, solves exactly this problem by letting a device on the local network resolve a friendly hostname such as \texttt{attendance.local} to whatever address the server currently holds, without needing a dedicated DNS server. This is the mechanism the ESP32 reader station uses to find the PC application, and it removes IP churn as a source of downtime entirely.

Communication between the reader station and the PC application follows standard HTTP semantics as defined in \textcite{fielding2014http}: the device issues \texttt{GET} requests to check server health and pull the current card list, and \texttt{POST} requests to submit taps, either individually or as a batch once a queue has built up. Using a well-understood, text-based protocol rather than a custom binary format keeps the system easy to test with ordinary tools during development and easy for future maintainers to extend.

\section{Software Stack for the Server-Side Application}

On the PC side, the application is built with Next.js, a React-based web framework documented in \textcite{vercel2026nextjs}, chosen for its ability to host both the API routes the ESP32 talks to and the dashboard staff use, within a single project. Data is persisted with Prisma, an object-relational mapper described in \textcite{prisma2026docs}, on top of SQLite \parencite{sqlite2026about}. SQLite was chosen deliberately over a client-server database: it stores the entire dataset in a single file on the PC's disk, needs no separate database server process to install or keep running, and is more than capable of handling the write volume a single-entrance attendance station generates. That combination keeps the "install and run" story for the PC application simple, which matters for a system that is meant to run unattended at a front desk rather than be managed by dedicated IT staff.

\section{Technology Adoption in the Local Context}

Low-cost microcontroller platforms are increasingly accessible in Zimbabwe specifically, not just in the literature at large. \textcite{techreviewafrica2026zimbabwe} reports that the Zimbabwean government is making IoT and robotics compulsory components of the school curriculum, reflecting a broader push toward local capacity to build rather than only consume this kind of technology. That context matters for a project like this one: the ESP32 and RC522 modules used here are inexpensive, widely available, and well within the budget of a small-to-medium organisation such as Infor Relay Systems (Pvt) Ltd, which is part of why this hardware combination was chosen over a more expensive commercial attendance terminal.

\section{Summary and Research Gap}

Taken together, the literature confirms two things. First, RFID is a well-established and appropriate technology for automating attendance capture, and a substantial body of recent student work has already proven the ESP32-plus-RC522 combination to be a workable, low-cost foundation \parencite{singh2022iot, yadav2025rfid, hutabarat2025implementation, turpo2025edurfid}. Second, almost none of that existing work addresses what happens when the network connection the system depends on is not reliable — a gap explicitly identified by \textcite{ishaq2023iot} across the wider field, and one that offline-first design patterns from the broader IoT literature can close \parencite{microsoft2026offline, topic2026offline, metadesk2026offline}. This project's contribution is not a new way of reading an RFID card; it is applying an offline-first architecture, mDNS-based discovery, and a dual-LED status indicator to that already well-proven RFID pattern, so that the resulting system keeps working during exactly the kind of network interruption a factory floor is likely to produce.
